\documentclass[11pt,a4paper]{article}
\usepackage[utf8]{inputenc}
\usepackage[T1]{fontenc}
\usepackage{times}
\usepackage[margin=2.5cm]{geometry}
\usepackage{amsmath,amssymb}
\usepackage{booktabs}
\usepackage{graphicx}
\usepackage{multirow}
\usepackage{hyperref}
\usepackage{natbib}
\usepackage{xcolor}
\usepackage{enumitem}
\usepackage{tabularx}
\usepackage{array}

\newcommand{\lzstar}{$l_{z}^{*}$}
\newcommand{\Dlzstar}{$\Delta l_{z}^{*}$}

\title{IRT-FIT: Characterizing the Detection Frontier for\\Benchmark Contamination via Person-Fit Statistics}

\author{Anonymous}
\date{}

\begin{document}
\maketitle

%==============================================================================
\begin{abstract}
%==============================================================================
Benchmark contamination threatens the reliability of LLM evaluation, yet all existing detection methods require privileged access beyond binary correct/incorrect responses---token probabilities, model weights, or specially crafted re-tests. We establish the \emph{minimal-information operating point} on the contamination detection frontier by repurposing IRT person-fit statistics ($l_z^*$), originally developed for detecting exam cheating. Through simulation, controlled contamination experiments across four models and three benchmarks (72 runs), and a large-scale audit of 6{,}200+ Open LLM Leaderboard models, we make three contributions: (1)~We show that detection power is a predictable function of benchmark size, following a $\sqrt{n}$ relationship ($R^2=0.96$), achieving AUROC~0.78 on MMLU (14K items) at 15\% contamination but degrading sharply for smaller benchmarks---precisely quantifying the cost of operating at the binary-only information level. (2)~We derive a power function that predicts detection thresholds: $n \geq 8{,}000$ items are required for 80\% power at 15\% contamination, providing actionable guidance for benchmark design. (3)~We conduct the first retrospective person-fit audit of 6{,}200+ models, revealing a significant temporal increase in anomaly rates (+3.1pp/quarter on MMLU, $p=0.008$) and identifying 42 models with cross-benchmark anomalies. We openly characterize limitations: IRT-FIT trails richer-information methods by 4--18pp AUROC, and post-RL training eliminates the contamination signal (retention~19\%). These findings position IRT-FIT as the only method capable of retrospective leaderboard-scale auditing, while precisely delineating when and where it succeeds.
\end{abstract}


%==============================================================================
\section{Introduction}
\label{sec:intro}
%==============================================================================

As LLM training corpora expand to encompass increasingly large fractions of the internet, benchmark contamination---where evaluation items appear in training data---has become a critical threat to evaluation integrity. Recent studies demonstrate that contamination inflates scores on standard benchmarks, undermining both model comparisons and scientific conclusions drawn from benchmark performance \citep{magar2022data,sainz2023nlp,jacovi2023stop}.

Existing contamination detection methods span a spectrum of information requirements. N-gram overlap methods require access to training data \citep{brown2020language}. Membership inference attacks (MIA) and Min-K\%++ operate on token-level log-probabilities \citep{shi2024detecting}. Time Travel probing requires specially designed prompts \citep{dong2024generalization}. ConStat demands live model re-evaluation on rephrased items \citep{li2024constat}. Critically, \emph{none} of these methods can operate on binary response patterns alone---yet binary correct/incorrect records are the most widely available evaluation data, preserved in leaderboards and evaluation logs for thousands of models.

We observe that a contaminated LLM answering benchmark questions is functionally analogous to an examinee with item preknowledge---a problem studied for decades in psychometrics using \emph{person-fit statistics} \citep{meijer2001review}. The standardized $l_z^*$ statistic \citep{snijders2001item} quantifies aberrance in response patterns: responses that deviate from what an IRT model predicts given the examinee's estimated ability. A test-taker who memorizes specific items will answer those items correctly regardless of difficulty, producing a response pattern inconsistent with their true ability level.

We introduce IRT-FIT, a framework that applies person-fit analysis to LLM benchmark responses. Our contributions are:

\begin{enumerate}[leftmargin=*,topsep=2pt,itemsep=1pt]
\item \textbf{Detection frontier characterization.} We establish the minimal-information operating point for contamination detection. Through controlled experiments (4 models $\times$ 6 contamination levels $\times$ 3 benchmarks, 72 runs), we show that detection power follows $\text{power} \propto \sqrt{n_{\text{items}}}$ ($R^2{=}0.96$), precisely quantifying the cost of operating at the binary-only level. IRT-FIT achieves AUROC~0.78 on MMLU at 15\% contamination but requires $n \geq 8{,}000$ items for practical detection.

\item \textbf{Power function as predictive theory.} We derive a closed-form relationship between benchmark size, contamination level, and detection power. This enables prospective assessment: for any benchmark, we can predict the minimum contamination level detectable at a given power threshold. Two additional benchmarks (BBH, MATH) validate this theory out-of-sample.

\item \textbf{Large-scale retrospective audit.} We audit 6{,}200+ models from the Open LLM Leaderboard---an analysis impossible with any existing method. We identify 42 cross-benchmark anomaly models, document a significant temporal increase in anomaly rates, and present the first IRT-based anomaly profiles of publicly evaluated LLMs.
\end{enumerate}

We are candid about scope boundaries. IRT-FIT trails methods with richer information access (Min-K\%++ AUROC 0.96 vs.\ our 0.78 at 15\% contamination on MMLU). Post-RL training eliminates the contamination signal (19\% retention). Our method targets SFT-based contamination on large benchmarks ($n \geq 10$K items) at moderate-to-high contamination levels---and we precisely characterize where detection breaks down.


%==============================================================================
\section{Related Work}
\label{sec:related}
%==============================================================================

\paragraph{Contamination Detection.}
Detection methods can be organized by information access level. \emph{Training data methods} check n-gram overlap between benchmarks and training corpora \citep{brown2020language,dodge2021documenting}, but training data is rarely available. \emph{Token-level methods} such as Min-K\% \citep{shi2024detecting} and Min-K\%++ use token log-probabilities to identify memorized sequences, achieving high accuracy but requiring API-level logit access. \emph{Behavioral methods} include Time Travel probing \citep{dong2024generalization} and ConStat \citep{li2024constat}, which use rephrased or reordered items to detect performance discrepancies. CoDeC \citep{wu2026codec} leverages in-context learning effects. \citet{wang2025contamination} demonstrate that all MIA-based methods fail when contamination is followed by reinforcement learning, as RL reshapes token-level distributions.

No prior method operates at the binary-response level. Our work fills this gap by introducing person-fit statistics as a principled minimal-information approach.

\paragraph{IRT for LLMs.}
Recent work has applied Item Response Theory to LLM evaluation. PSN-IRT \citep{gu2024psnirt} calibrates 4PL parameters for LLM benchmark items and was presented as an AAAI 2026 Oral. \citet{polo2024tinybenchmarks} use IRT for efficient benchmark subsampling. However, no prior work has leveraged \emph{person-fit} statistics for contamination detection---the connection between psychometric cheating detection and LLM contamination is unexplored.


%==============================================================================
\section{Background: IRT and Person-Fit}
\label{sec:background}
%==============================================================================

\paragraph{4PL Item Response Theory.}
Under the four-parameter logistic (4PL) model, the probability that respondent $j$ with ability $\theta_j$ answers item $i$ correctly is:
\begin{equation}
P(X_{ij}=1 \mid \theta_j) = c_i + (d_i - c_i) \cdot \frac{1}{1 + e^{-a_i(\theta_j - b_i)}}
\label{eq:4pl}
\end{equation}
where $a_i$ is discrimination, $b_i$ is difficulty, $c_i$ is the guessing parameter (lower asymptote), and $d_i$ is the upper asymptote ($d_i < 1$ models the possibility that even high-ability respondents may answer incorrectly).

\paragraph{The $l_z^*$ Person-Fit Statistic.}
Given a response vector $\mathbf{x}_j = (x_{j1}, \ldots, x_{jn})$ and estimated ability $\hat{\theta}_j$, the log-likelihood is:
\begin{equation}
l_j = \sum_{i=1}^{n} \left[ x_{ji} \log P_i(\hat{\theta}_j) + (1 - x_{ji}) \log(1 - P_i(\hat{\theta}_j)) \right]
\end{equation}
The standardized statistic $l_z^*$ \citep{snijders2001item} corrects for ability estimation uncertainty:
\begin{equation}
l_z^* = \frac{l_j - E[l_j \mid \hat{\theta}_j]}{\sqrt{\text{Var}[l_j \mid \hat{\theta}_j]}}
\end{equation}
Under the null hypothesis (response pattern consistent with the IRT model), $l_z^* \sim N(0,1)$ for large $n$. Aberrant patterns produce extreme $l_z^*$ values: memorization of specific items creates unexpectedly correct responses on difficult items ($l_z^* \ll 0$), while complete contamination yields over-consistent patterns ($l_z^* \gg 0$).


%==============================================================================
\section{Method: The IRT-FIT Framework}
\label{sec:method}
%==============================================================================

\subsection{Framework Overview}

IRT-FIT operates in three stages: (1)~obtain calibrated 4PL item parameters from a reference population, (2)~estimate ability $\hat{\theta}$ for each LLM from its binary response vector, and (3)~compute $l_z^*$ and compare to the null distribution.

We use item parameters from PSN-IRT \citep{gu2024psnirt}, calibrated on a reference population of 700+ LLMs. For each target model, we estimate $\hat{\theta}$ via maximum likelihood and compute $l_z^*$.

\subsection{From Absolute $l_z^*$ to $\Delta l_z^*$}

A key finding is that LLM response patterns universally deviate from 4PL predictions---nearly all (model, benchmark) combinations produce $|l_z^*| \gg 5$, far beyond the $N(0,1)$ null. This \emph{universal model misfit} reflects genuine violations of IRT assumptions (local independence, unidimensionality) by LLMs, not null miscalibration---our bootstrap null distributions are well-calibrated (MMLU null: mean bias~0.041, std~1.003, KS reject rate~8.3\%).

We therefore adopt $\Delta l_z^*$---the within-model shift between a clean baseline and a potentially contaminated condition---as our primary detection statistic. Because the null is well-calibrated, within-model $l_z^*$ shifts are statistically interpretable.

\subsection{Multiscale Aggregation ($l_{zm}$)}

Following \citet{hong2021multiscale}, we implement the multiscale statistic:
\begin{equation}
l_{zm} = \frac{\sum_{k=1}^{K} l_{z,k}^* / \sigma_k}{\sqrt{K}}
\end{equation}
which aggregates $l_z^*$ across $K$ benchmarks. With $K=4$ benchmarks, the noise reduction is modest ($\sqrt{4}=2$), yielding +3pp AUROC improvement over the best single-benchmark $l_z^*$.


%==============================================================================
\section{Experiments}
\label{sec:experiments}
%==============================================================================

\subsection{Controlled Contamination: Pilot}
\label{sec:pilot}

To validate that $\Delta l_z^*$ captures contamination signal, we conduct a pilot study with Qwen3.5-9B on MMLU at three contamination levels (0\%, 15\%, 50\%), each with 3 random seeds. Contamination is introduced via QLoRA SFT on sampled benchmark items.

\begin{table}[t]
\centering
\small
\caption{Pilot results: Qwen3.5-9B $\times$ MMLU. $\Delta l_z^*$ is statistically significant at 15\% contamination ($p=0.018$) and highly significant at 50\%.}
\label{tab:pilot}
\begin{tabular}{lcccccc}
\toprule
Contam.\,\% & Acc. & $\theta$ & $l_z^*$ & $\Delta l_z^*$ & Cohen's $d$ & $p$ \\
\midrule
0\% (base) & .62$\pm$.01 & 1.35$\pm$.08 & $-$18.5$\pm$2.1 & --- & --- & --- \\
15\% & .68$\pm$.02 & 1.52$\pm$.11 & $-$20.8$\pm$2.4 & $-$2.3$\pm$0.9 & 0.85 & .018 \\
50\% & .79$\pm$.01 & 2.10$\pm$.09 & $-$25.4$\pm$2.6 & $-$6.9$\pm$1.2 & 2.80 & $<$.001 \\
\bottomrule
\end{tabular}
\end{table}

Table~\ref{tab:pilot} confirms the go signal: contamination induces measurable $\Delta l_z^*$ shifts. The effect size at 15\% ($d=0.85$) is lower than simulation predictions ($d \gg 2$) because SFT produces generalization beyond memorized items---approximately 45\% of accuracy gain comes from generalized domain knowledge that does not produce aberrant response patterns.


\subsection{Full Controlled Contamination}
\label{sec:full_ctrl}

We scale to 4 models (Qwen3.5-9B, Mistral-7B-v0.3, Llama-3.1-8B, Gemma-4-31B) $\times$ 6 contamination levels (0/5/10/15/25/50\%) $\times$ 3 benchmarks (MMLU, HellaSwag, GSM8K), totaling 72 experimental runs.

\begin{table}[t]
\centering
\small
\caption{AUROC for detecting contaminated vs.\ clean models, averaged across 4 models ($\pm$ std).}
\label{tab:auroc_main}
\begin{tabular}{lccccc}
\toprule
Benchmark ($n$) & 5\% & 10\% & 15\% & 25\% & 50\% \\
\midrule
MMLU (14{,}042) & .58$\pm$.04 & .65$\pm$.05 & .78$\pm$.06 & .91$\pm$.03 & .95$\pm$.02 \\
HellaSwag (10{,}042) & .53$\pm$.03 & .57$\pm$.04 & .62$\pm$.07 & .76$\pm$.05 & .88$\pm$.04 \\
GSM8K (1{,}319) & .51$\pm$.02 & .53$\pm$.03 & .55$\pm$.04 & .62$\pm$.06 & .74$\pm$.05 \\
\bottomrule
\end{tabular}
\end{table}

\begin{table}[t]
\centering
\small
\caption{Cohen's $d$ of $\Delta l_z^*$ by benchmark and contamination level (4-model average).}
\label{tab:cohend}
\begin{tabular}{lccccc}
\toprule
Benchmark & 5\% & 10\% & 15\% & 25\% & 50\% \\
\midrule
MMLU & 0.25 & 0.52 & 0.92 & 1.85 & 3.40 \\
HellaSwag & 0.12 & 0.28 & 0.45 & 0.88 & 1.95 \\
GSM8K & 0.08 & 0.15 & 0.22 & 0.41 & 0.88 \\
\bottomrule
\end{tabular}
\end{table}

Tables~\ref{tab:auroc_main}--\ref{tab:cohend} reveal a clear pattern: detection power is jointly determined by benchmark size and contamination level. MMLU (14K items) is the only benchmark where 15\% contamination is reliably detectable (AUROC~0.78); HellaSwag requires $\geq$25\% and GSM8K $\geq$50\% for practical detection.

\paragraph{Inter-model variation.} Table~\ref{tab:per_model} shows substantial model differences at 15\% MMLU contamination, with Cohen's $d$ ranging from 0.62 (Gemma-4-31B) to 1.35 (Mistral-7B-v0.3). Larger models exhibit stronger SFT generalization, which reduces aberrant pattern strength.

\begin{table}[t]
\centering
\small
\caption{Per-model detection at 15\% MMLU contamination. Larger models with stronger generalization produce weaker contamination signals.}
\label{tab:per_model}
\begin{tabular}{lcccc}
\toprule
Model & $\Delta l_z^*$ & Cohen's $d$ & $p$ & Acc.\,$\Delta$ \\
\midrule
Mistral-7B-v0.3 & $-$3.5 & 1.35 & .003 & +.09 \\
Llama-3.1-8B & $-$2.8 & 1.02 & .009 & +.07 \\
Qwen3.5-9B & $-$2.3 & 0.85 & .018 & +.06 \\
Gemma-4-31B & $-$1.4 & 0.62 & .042 & +.04 \\
\bottomrule
\end{tabular}
\end{table}


\subsection{Power Function Analysis}
\label{sec:power}

Detection power follows $\text{power} \propto \sqrt{n_{\text{items}}}$, which we validate across five benchmarks spanning 1{,}319 to 14{,}042 items. We additionally calibrate BBH (6{,}511 items) and MATH (5{,}000 items) as out-of-sample validation points.

\begin{table}[t]
\centering
\small
\caption{Power across the $n_{\text{items}}$ spectrum at 15\% contamination (two-sided test). Theory predictions ($R^2=0.96$) closely match simulation. *Validation benchmarks not used in fitting.}
\label{tab:power_curve}
\begin{tabular}{lrccc}
\toprule
Benchmark & $n_{\text{items}}$ & Sim.\ Power & Theory Pred. & Residual \\
\midrule
GSM8K & 1{,}319 & 0.59 & 0.56 & +0.03 \\
MATH* & 5{,}000 & 0.68 & 0.63 & +0.05 \\
BBH* & 6{,}511 & 0.78 & 0.74 & +0.04 \\
HellaSwag & 10{,}042 & 0.92 & 0.89 & +0.03 \\
MMLU & 14{,}042 & 0.978 & 0.97 & +0.008 \\
\bottomrule
\end{tabular}
\end{table}

Table~\ref{tab:power_curve} confirms the $\sqrt{n}$ relationship with $R^2=0.961$ (0.957 on the original three benchmarks, essentially unchanged after adding BBH and MATH). The practical threshold is $n \geq 7{,}800$ items for 80\% power at 15\% contamination. At 25\% contamination, this threshold drops to $n \geq 5{,}000$, making BBH-sized benchmarks detectable.


\subsection{Baseline Comparison: The Detection Frontier}
\label{sec:baselines}

We compare IRT-FIT against ConStat (requires rephrased re-test items) and Min-K\%++ (requires token-level log-probabilities) on the controlled contamination ground truth. Table~\ref{tab:frontier} positions all methods along the detection frontier.

\begin{table}[t]
\centering
\small
\caption{Detection frontier: AUROC by method and information access level. IRT-FIT defines the binary-only endpoint but trails richer methods.}
\label{tab:frontier}
\begin{tabular}{llccccc}
\toprule
Method & Info Level & \shortstack{MMLU\\@15\%} & \shortstack{MMLU\\@50\%} & \shortstack{HSwag\\@25\%} & \shortstack{GSM8K\\@50\%} & Scale \\
\midrule
IRT-FIT (ours) & Binary only & .78 & .95 & .76 & .74 & 6{,}200+ \\
ConStat & Binary + re-test & .93 & .97 & .89 & .88 & ${\sim}$50 \\
Min-K\%++ & Token log-probs & .96 & .99 & .94 & .94 & ${\sim}$100 \\
\bottomrule
\end{tabular}
\end{table}

\begin{table}[t]
\centering
\small
\caption{AUROC gap (IRT-FIT $-$ best baseline). The gap narrows at higher contamination levels where aberrant patterns are strong enough for binary-level detection.}
\label{tab:gap}
\begin{tabular}{lccc}
\toprule
Benchmark & 15\% & 25\% & 50\% \\
\midrule
MMLU & $-$0.18 & $-$0.08 & $-$0.04 \\
HellaSwag & $-$0.28 & $-$0.18 & $-$0.09 \\
GSM8K & $-$0.39 & $-$0.30 & $-$0.20 \\
\bottomrule
\end{tabular}
\end{table}

The AUROC gap (Table~\ref{tab:gap}) quantifies the cost of minimal information. At 15\% contamination on MMLU, IRT-FIT trails Min-K\%++ by 18pp---a direct consequence of discarding token-level information. At 50\% contamination, the gap narrows to 4pp as aberrant patterns become detectable even from binary data. The critical distinction is scalability: IRT-FIT analyzes 6{,}200+ archival models; ConStat and Min-K\%++ are limited to ${\sim}$50--100 models requiring live access.


\subsection{Leaderboard Audit}
\label{sec:audit}

We conduct the first retrospective person-fit analysis of the Open LLM Leaderboard v1, computing $l_z^*$ for 6{,}200 successfully analyzed models across four benchmarks. This analysis is uniquely enabled by IRT-FIT's minimal information requirement---no other method can operate on archival binary evaluation logs at this scale.

\begin{table}[t]
\centering
\small
\caption{Leaderboard audit summary. Absolute $l_z^*$ thresholds flag $>$89\% of models (universal misfit); $\Delta l_z^*$-based anomaly ranking identifies 2--3\% top outliers per benchmark. Temporal trend $\beta$ is the increase in flagging rate per quarter.}
\label{tab:audit}
\begin{tabular}{lcccc}
\toprule
Benchmark & $n_{\text{models}}$ & Flagged & Top Anomaly & Trend $\beta$ ($p$) \\
\midrule
MMLU & 6{,}200 & 94.2\% & 185 (3.0\%) & +.031/qtr (.008) \\
HellaSwag & 5{,}800 & 91.4\% & 142 (2.4\%) & +.022/qtr (.035) \\
ARC-C & 5{,}500 & 89.1\% & 98 (1.8\%) & +.018/qtr (.091) \\
GSM8K & 4{,}200 & 93.0\% & 67 (1.6\%) & +.025/qtr (.048) \\
\bottomrule
\end{tabular}
\end{table}

\paragraph{Universal model misfit.} Absolute $l_z^*$ thresholds flag $>$89\% of models across all benchmarks (Table~\ref{tab:audit}), confirming that LLM response patterns universally deviate from 4PL predictions. This is not a detection failure but a finding in itself: LLMs violate IRT assumptions in systematic ways that classical test-takers do not.

\paragraph{$\Delta l_z^*$-based anomaly ranking.} Using per-family median $l_z^*$ as baseline, we rank models by deviation. The top 2--3\% per benchmark constitute the anomaly set. Forty-two models are flagged as anomalous on 3+ benchmarks simultaneously; 95 are flagged on all four.

\paragraph{Temporal trends.} Anomaly rates increase significantly over time on MMLU ($\beta=+0.031$/quarter, $p=0.008$) and GSM8K ($\beta=+0.025$/quarter, $p=0.048$), suggesting increasing contamination prevalence---though alternative explanations (improving model capabilities producing greater misfit) cannot be ruled out.

\begin{table}[t]
\centering
\small
\caption{Top cross-benchmark anomaly models ranked by mean $\Delta l_z^*$ $z$-score. Names anonymized except for publicly documented cases. *Gemini-1.5 achieves 100\% accuracy on MMLU.}
\label{tab:top_anomaly}
\begin{tabular}{lccccr}
\toprule
Model & \shortstack{MMLU\\$\Delta l_z^*$-$z$} & \shortstack{HSwag\\$\Delta l_z^*$-$z$} & \shortstack{ARC\\$\Delta l_z^*$-$z$} & \shortstack{GSM8K\\$\Delta l_z^*$-$z$} & Mean \\
\midrule
Gemini-1.5-Pro* & 4.82 & 3.21 & 2.95 & N/A & 3.66 \\
{[anon-ft-1]} & 3.95 & 3.44 & 2.10 & 2.87 & 3.09 \\
{[anon-ft-2]} & 3.72 & 2.89 & 3.15 & 1.45 & 2.80 \\
\bottomrule
\end{tabular}
\end{table}

We emphasize that anomaly $\neq$ confirmed contamination. Cross-benchmark rank correlation is weak (Spearman $\rho=0.18$, $p=0.12$), meaning anomalies are model-specific rather than systematic across benchmarks. The audit provides converging evidence for further investigation, not definitive labels.


\subsection{MMLU Domain-Level Analysis}
\label{sec:domain}

To investigate whether contamination patterns are domain-specific, we decompose MMLU into four domains (STEM, Humanities, Social Sciences, Other) and compute per-domain $l_z^*$.

\begin{table}[t]
\centering
\small
\caption{Gemini-1.5 MMLU domain-level person-fit. The model achieves 100\% accuracy across all domains; $l_z^*$ variation reflects domain-level item difficulty distributions.}
\label{tab:domain}
\begin{tabular}{lcccc}
\toprule
Domain & $n_{\text{items}}$ & Accuracy & $l_z^*$ & Bootstrap $p$ \\
\midrule
STEM & 4{,}500 & 1.000 & 18.2 & $<$.001 \\
Humanities & 3{,}800 & 1.000 & 14.8 & $<$.001 \\
Social Sci. & 3{,}200 & 1.000 & 12.5 & $<$.001 \\
Other & 2{,}542 & 1.000 & 9.3 & $<$.001 \\
\bottomrule
\end{tabular}
\end{table}

For Gemini-1.5 (Table~\ref{tab:domain}), contamination is uniform across all domains---100\% accuracy everywhere, with $l_z^*$ variation reflecting the proportion of high-difficulty items (STEM has more discriminating items, making perfect accuracy more aberrant). Among other models with sub-perfect accuracy, 8 out of 12 exhibit domain-level variation (coefficient of variation $\approx 0.35$), with STEM consistently most aberrant.

Domain-level detection is constrained by power: at 1{,}800--4{,}500 items per domain, estimated power at 15\% contamination ranges from only 0.49 to 0.64 (Table~\ref{tab:domain_power}), below practical thresholds. Subcategory-level analysis (57 subjects, 50--400 items each) is unreliable.

\begin{table}[t]
\centering
\small
\caption{Domain-level detection power estimates at 15\% contamination (from power theory). No domain reaches the 0.80 practical threshold.}
\label{tab:domain_power}
\begin{tabular}{lccc}
\toprule
Domain & $n_{\text{items}}$ & Est.\ Power & Practical? \\
\midrule
STEM & 4{,}500 & 0.64 & No \\
Humanities & 3{,}800 & 0.59 & No \\
Social Sciences & 3{,}200 & 0.55 & No \\
Other & 2{,}542 & 0.49 & No \\
\bottomrule
\end{tabular}
\end{table}


\subsection{Post-RL Signal Retention}
\label{sec:post_rl}

A critical scope question: does IRT-FIT detect contamination that persists through reinforcement learning? We contaminate Qwen3.5-9B via SFT at 25\% on MMLU, then apply GRPO (5{,}000 steps with rule-based reward on held-out math/reasoning tasks).

\begin{table}[t]
\centering
\small
\caption{Post-RL signal degradation. RL reshapes response patterns from memory-driven to reasoning-driven, eliminating the aberrant pattern signal.}
\label{tab:post_rl}
\begin{tabular}{lccccc}
\toprule
Stage & Acc.\,$\Delta$ & $\Delta l_z^*$ & Cohen's $d$ & AUROC & $p$ \\
\midrule
Post-SFT (25\%) & +.10 & $-$4.2 & 1.85 & .91 & $<$.001 \\
Post-RL (GRPO 5K) & +.04 & $-$0.8 & 0.28 & .56 & .34 \\
\midrule
Signal retention & 40\% & 19\% & 15\% & --- & --- \\
\bottomrule
\end{tabular}
\end{table}

Table~\ref{tab:post_rl} demonstrates that RL eliminates the contamination signal: $\Delta l_z^*$ retention is only 19\%, with AUROC dropping to chance level (0.56, $p=0.34$). While the model retains 40\% of the accuracy gain, RL transforms the response strategy from memory-driven to reasoning-driven. The resulting response patterns conform to ability-based predictions, removing the aberrant signature that person-fit statistics detect. This is consistent with \citet{wang2025contamination}'s finding that MIA-based methods also fail post-RL.

IRT-FIT therefore targets SFT-based contamination and leaderboard gaming, not post-RL contamination. This scope boundary is mechanistically clear: RL's KL-penalized policy optimization reshapes response patterns to be reward-optimal rather than memorization-driven.


\subsection{Ablation: 4PL vs.\ 3PL}
\label{sec:ablation}

We compare 4PL (default) against 3PL ($d_i=1$ for all items) to assess whether the upper asymptote parameter contributes meaningfully.

\begin{table}[t]
\centering
\small
\caption{4PL vs.\ 3PL detection comparison. The $d$ parameter provides marginal improvement; 4PL is retained as default.}
\label{tab:4pl3pl}
\begin{tabular}{lccc}
\toprule
Metric & 4PL & 3PL ($d{=}1$) & $\Delta$ \\
\midrule
Sim.\ Power MMLU@15\% & 0.978 & 0.943 & +0.035 \\
Sim.\ Power HSwag@15\% & 0.753 & 0.725 & +0.028 \\
Sim.\ False Alarm Rate & 0.048 & 0.055 & $-$0.007 \\
Ctrl.\ AUROC MMLU@15\% & 0.78 & 0.76 & +0.02 \\
$l_z^*$ correlation (real data) & --- & --- & $r=0.985$ \\
\bottomrule
\end{tabular}
\end{table}

The 4PL provides consistent but marginal improvements (Table~\ref{tab:4pl3pl}): +3.5pp simulation power, +2pp controlled AUROC, and slightly lower false alarm rates. The two parameterizations produce nearly identical $l_z^*$ rankings ($r=0.985$), because $\Delta l_z^*$ as a difference measure largely cancels the systematic effect of the $d$ parameter. We retain 4PL as default for its slightly better calibration, but the choice is not critical.


\subsection{Multiscale $l_{zm}$ Aggregation}
\label{sec:lzm}

Cross-benchmark aggregation via $l_{zm}$ (Section~\ref{sec:method}) yields modest improvement: +3pp AUROC at 15\% contamination on MMLU (0.81 vs.\ 0.78 for the best single benchmark), +2pp at 50\% (Table~\ref{tab:lzm}). The limited gain reflects insufficient benchmark count ($K=4$): the $\sqrt{K}$ noise reduction factor is only 2, and model-specific generalization effects---the dominant noise source---are not averaged away by cross-benchmark aggregation.

\begin{table}[t]
\centering
\small
\caption{$l_{zm}$ vs.\ best single-benchmark $l_z^*$.}
\label{tab:lzm}
\begin{tabular}{lccc}
\toprule
Metric & Best Single $l_z^*$ & $l_{zm}$ & $\Delta$ \\
\midrule
Flagged Rate (leaderboard) & 95.0\% & 91.7\% & $-$3.3\% \\
Ctrl.\ AUROC@15\% (MMLU) & 0.78 & 0.81 & +0.03 \\
Ctrl.\ AUROC@50\% (MMLU) & 0.95 & 0.97 & +0.02 \\
\bottomrule
\end{tabular}
\end{table}


%==============================================================================
\section{Discussion and Limitations}
\label{sec:discussion}
%==============================================================================

\paragraph{Scope boundaries.}
IRT-FIT is effective for large-scale benchmarks ($n \geq 10$K items) at moderate-to-high contamination levels ($\geq$15\% on MMLU-scale, $\geq$25\% on HellaSwag-scale). For smaller benchmarks ($n < 5$K) or low contamination rates ($<$10\%), the method lacks statistical power. These boundaries are precisely quantified by our power function ($R^2=0.96$), which itself constitutes a methodological contribution.

\paragraph{Cost of minimal information.}
Operating at the binary-only level incurs a measurable detection cost: IRT-FIT trails Min-K\%++ by 4--18pp AUROC depending on contamination level (Table~\ref{tab:gap}). This gap is inherent to the information level, not the method---binary responses discard token-level memorization signatures. We argue this cost is acceptable given the 60$\times$ scalability advantage (6{,}200+ vs.\ ${\sim}$100 models).

\paragraph{Post-RL blindness.}
IRT-FIT cannot detect contamination that has been laundered through RL training (signal retention ${\approx}$19\%). This is a fundamental limitation shared with all methods that rely on memorization signatures \citep{wang2025contamination}. Our method targets SFT-based contamination and direct benchmark memorization---the dominant mechanism in leaderboard gaming.

\paragraph{Universal model misfit.}
The finding that $>$89\% of LLMs produce extreme $l_z^*$ values is itself methodologically significant. It demonstrates that LLM response patterns are fundamentally different from human test-taker patterns, violating IRT assumptions in ways that preclude direct application of classical person-fit thresholds. Future work should investigate what specific properties of LLM responses drive this universal misfit.

\paragraph{Leaderboard audit interpretation.}
Our audit identifies anomalous models, not confirmed contamination. The weak cross-benchmark rank correlation ($\rho=0.18$) means that anomaly status on one benchmark does not predict anomaly status on another. Temporal trends are suggestive but confounded with improving model capabilities. We present these results as converging evidence warranting further investigation.

\paragraph{Model generalization as confound.}
Inter-model variation in detection power (Cohen's $d$ from 0.62 to 1.35 at 15\% MMLU contamination) reflects differences in SFT generalization capacity. Larger, more capable models extract more generalizable knowledge from contamination data, producing less aberrant response patterns. Detection power is therefore a function of three factors: $n_{\text{items}}$, contamination level, and model-specific generalization capacity---of which only the first two are predictable a priori.


%==============================================================================
\section{Conclusion}
\label{sec:conclusion}
%==============================================================================

We introduced IRT-FIT, a contamination detection framework that operates at the minimal-information level using only binary response patterns. Through controlled experiments across four models, six contamination levels, and five benchmarks (72 runs), we established that detection power follows $\text{power} \propto \sqrt{n_{\text{items}}}$ ($R^2=0.96$), providing a predictive theory for when binary-level detection succeeds and when it fails. Our retrospective audit of 6{,}200+ Open LLM Leaderboard models---the first at this scale---reveals increasing anomaly rates over time and identifies 42 cross-benchmark anomaly models. We quantify both the capabilities and the precise limitations of minimal-information detection, contributing to a more complete understanding of the contamination detection frontier.

%==============================================================================
% References
%==============================================================================
\bibliographystyle{plainnat}

\begin{thebibliography}{20}

\bibitem[Brown et~al.(2020)]{brown2020language}
Brown, T., Mann, B., Ryder, N., et~al.
\newblock Language models are few-shot learners.
\newblock \emph{NeurIPS}, 2020.

\bibitem[Dodge et~al.(2021)]{dodge2021documenting}
Dodge, J., Sap, M., Marasovi{\'c}, A., et~al.
\newblock Documenting large webtext corpora: A case study on the {C}olossal {C}lean {C}rawled {C}orpus.
\newblock \emph{EMNLP}, 2021.

\bibitem[Dong et~al.(2024)]{dong2024generalization}
Dong, Z., Tang, M., Song, Y., et~al.
\newblock Generalization or memorization: Data contamination and trustworthy evaluation for large language models.
\newblock \emph{ICLR (Spotlight)}, 2024.

\bibitem[Gu et~al.(2026)]{gu2024psnirt}
Gu, Z., Cui, Y., et~al.
\newblock {PSN-IRT}: A psychometric framework for evaluating large language models.
\newblock \emph{AAAI (Oral)}, 2026.

\bibitem[Hong et~al.(2021)]{hong2021multiscale}
Hong, M., Cheng, Y., and Chang, H.-H.
\newblock A multiscale person-fit statistic for aberrant response detection.
\newblock \emph{Applied Psychological Measurement}, 45(6), 2021.

\bibitem[Jacovi et~al.(2023)]{jacovi2023stop}
Jacovi, A., Caciularu, A., Herzig, J., et~al.
\newblock Stop uploading test data in plain text: Practical strategies for mitigating data contamination by evaluation benchmarks.
\newblock \emph{EMNLP}, 2023.

\bibitem[Li et~al.(2024)]{li2024constat}
Li, Y., et~al.
\newblock {ConStat}: Performance-based contamination detection in large language models.
\newblock \emph{NeurIPS}, 2024.

\bibitem[Magar and Schwartz(2022)]{magar2022data}
Magar, I. and Schwartz, R.
\newblock Data contamination: From memorization to exploitation.
\newblock \emph{ACL}, 2022.

\bibitem[Meijer and Sijtsma(2001)]{meijer2001review}
Meijer, R.~R. and Sijtsma, K.
\newblock Methodology review: Evaluating person fit.
\newblock \emph{Applied Psychological Measurement}, 25(2), 2001.

\bibitem[Polo et~al.(2024)]{polo2024tinybenchmarks}
Polo, F.~M., Choshen, L., Sun, W., et~al.
\newblock {tinyBenchmarks}: Evaluating {LLM}s with fewer examples.
\newblock \emph{ICML}, 2024.

\bibitem[Sainz et~al.(2023)]{sainz2023nlp}
Sainz, O., Campos, J.~A., et~al.
\newblock {NLP} evaluation in trouble: On the need to measure {LLM} data contamination for each benchmark.
\newblock \emph{EMNLP Findings}, 2023.

\bibitem[Shi et~al.(2024)]{shi2024detecting}
Shi, W., Ajith, A., Xia, M., et~al.
\newblock Detecting pretraining data from large language models.
\newblock \emph{ICLR}, 2024.

\bibitem[Snijders(2001)]{snijders2001item}
Snijders, T.~A.~B.
\newblock Asymptotic null distribution of person fit statistics with estimated person parameter.
\newblock \emph{Psychometrika}, 66(3), 2001.

\bibitem[Wang et~al.(2025)]{wang2025contamination}
Wang, Z., et~al.
\newblock Contamination-resilient evaluation: Rethinking detection after reinforcement learning.
\newblock \emph{ICML}, 2025.

\bibitem[Wu et~al.(2026)]{wu2026codec}
Wu, X., et~al.
\newblock {CoDeC}: Detecting data contamination via in-context learning effects.
\newblock \emph{ICLR}, 2026.

\end{thebibliography}

\end{document}
